\Subhead{5.2.2}{Тропосферные радиоволны. Рассеяние и поглощение. Замирания. Тропосферный волновод}

Тропосферный показатель близок к единице, но меняется с температурой, давлением и влажностью. Плавный градиент искривляет лучи, мелкие флуктуации рассеивают; молекулярные резонансы, облака и осадки поглощают/рассеивают СВЧ с выраженной частотной зависимостью. Поэтому свободнопространственную формулу ослабления нельзя без поправок применять на любой частоте и дальности.

Несколько путей с разными амплитудами и задержками создают
\begin{equation}
 E(t)=\sum_j a_js(t-\tau_j)\e^{\ii\phi_j}.
\end{equation}
Изменение фаз из-за движения и состояния среды приводит к \emph{замираниям}. При большом числе случайных вкладов без доминирующего прямого луча комплексная амплитуда близка к гауссовской, огибающая — к Рэлеевской; сильный регулярный луч ведёт к райсовской модели. Разброс задержек создаёт частотно-селективные замирания; при полосе $B\Delta\tau\ll1$ канал приблизительно плоский.

Сильный отрицательный вертикальный градиент показателя может создать тропосферный канал. С учётом кривизны Земли используют модифицированную рефрактивность $M=N+10^6h/R_E$, $N=10^6(n-1)$, $h\ll R_E$. Участок $\dd M/\dd h<0$ допускает захват пологих лучей при подходящих границах; это сильнее обычной атмосферной рефракции. Разнесение антенн, частот или поляризаций уменьшает одновременное глубокое замирание, если соответствующие каналы достаточно независимы.
